\chapter{GroundOne: metodo e implementazione}
\label{sec:method}

% Source notes: Project/docs/thesis_notes.md (Ch. 3), Project/docs/dataset_plan.md §8, implementation_plan.md §4.

GroundOne è un modello linguistico compatto usato come System One: legge una sola volta il documento e la risposta e,
senza generare testo, emette un verdetto per ciascuna frase valutata della risposta, ricavandolo dai logit di poche
lettere di risposta. Il capitolo descrive la decisione che il modello prende, il prompt che raccoglie in un unico
forward pass tutte le decisioni relative a una risposta, l'obiettivo di addestramento, la calibrazione delle
probabilità in uscita, la regola che converte le probabilità nelle etichette restituite e il software che implementa il
tutto.

\section{Task e decisione}
\label{sec:method-task}

Per ogni frase di una risposta GroundOne sceglie fra tre classi, definite nella Sezione~\ref{sec:data-task}: rispetto
al documento, la frase è \emph{supportata}, \emph{contraddetta} o \emph{non menzionata}. Alle tre classi corrispondono
le lettere A, B e C, e il modello restituisce una probabilità per ciascuna. Il punteggio binario di groundedness e le
policy strict e lenient non sono domande distinte poste al modello, ma vengono calcolati a partire da queste tre
probabilità, cosicché un unico forward pass basta per ogni utilizzo. Interrogare il modello direttamente sulle policy
non porterebbe nulla: ciascuna è una funzione deterministica della classe a tre vie, per cui domande separate
moltiplicherebbero le decisioni senza aggiungere informazione e potrebbero fornire al modello obiettivi di
addestramento incoerenti per la stessa frase.

\section{Tutti i verdetti in un solo forward pass}
\label{sec:method-readout}

Il prompt contiene una breve istruzione, il documento e la risposta. Le frasi della risposta sono numerate, e ognuna di
quelle a cui è associata una decisione è seguita da un marker, la stringa ``\texttt{~Verdict:}''. Le frasi prive di
contenuto verificabile (Sezione~\ref{sec:data-construction}) restano nel testo come contesto, ma senza marker. Il
modello elabora l'intera sequenza una sola volta: sull'ultimo token di ciascun marker si leggono i logit delle lettere
A, B e C, normalizzati con una softmax ristretta a queste tre lettere. Poiché in un decoder ogni posizione vede
soltanto quelle precedenti, il verdetto sulla frase $i$ dipende dal documento e dalle frasi $1, \dots, i$, ed è
esattamente ciò che serve alla verifica, dato che i riferimenti, come i pronomi, puntano all'indietro. Quanto segue
l'ultimo marker non influisce su alcun verdetto, quindi la sequenza si chiude lì.

Due accorgimenti rendono il readout efficiente ed esatto. Anzitutto il prompt viene costruito tokenizzando
separatamente le sue parti e concatenandole, così che la posizione di ogni marker sia nota con esattezza e resti
identica in addestramento e in inferenza. In secondo luogo, la proiezione di output sul vocabolario si applica solo
nelle posizioni dei marker: il corpo del transformer elabora l'intera sequenza, ma la grande matrice del vocabolario
viene moltiplicata soltanto per i pochi stati nascosti effettivamente letti.

Questa scelta determina il costo di una verifica. Indicando con $S$ la lunghezza complessiva di documento e risposta,
con $K$ il numero di decisioni e con $m$ la lunghezza in token di un marker, il System One richiede un solo prefill di
$S + K m$ token. Un giudice generativo deve invece produrre anche i propri verdetti in fase di decoding: con $n$ token
per verdetto e un tempo di decoding per token $t_\mathrm{dec}$, il suo costo è all'incirca
$\mathrm{prefill}(S) + K\, n\, t_\mathrm{dec}$. Dato che $m$ vale tre token e che sull'hardware di riferimento il
decoding è, a parità di token, molto più lento del prefill, il divario si allarga al crescere di $K$. Sull'M4 usato in
questa tesi il prefill costa circa 0.37\,ms per token con il modello 0.5B e 2.2\,ms con il 3B, mentre il decoding costa
5.9 e 31.8\,ms per token generato: un token in decoding costa quanto quattordici-sedici token in prefill. Le due fasi,
inoltre, sollecitano l'hardware in modo diverso. Il prefill elabora molti token insieme ed è limitato dalla capacità di
calcolo, mentre il decoding produce un token alla volta ed è limitato dalla banda necessaria a leggere i
pesi~\cite{pope2023scaling,patel2024splitwise}. La differenza riemerge nelle misure di energia del
Capitolo~\ref{sec:results}.

Una seconda soluzione legge una sola volta il prefisso condiviso e vi accoda poi una breve domanda per ogni frase,
valutando le domande in batch. Richiede però una KV cache replicata $K$ volte, mentre la soluzione a marker si
accontenta di un solo passaggio. È mantenuta come alternativa di riserva ed è il formato naturale per un modello non
addestrato, che comprende una domanda esplicita ma non ha motivo di interpretare un marker come tale; per questo è il
formato usato dalla baseline non addestrata (Sezione~\ref{sec:setup-baselines}).

\paragraph{Un classificatore, non un generatore più veloce.} Il readout restringe la decisione a tre lettere, e la
restrizione non è un dettaglio. Dopo il marker, il modello sottoposto a fine-tuning non scriverebbe la lettera che viene
letta: il suo token successivo più probabile è la stessa lettera preceduta da uno spazio, e le tre lettere del readout
raccolgono insieme appena il 4--5\,\% circa della massa di probabilità in corrispondenza dei marker. Addestrare con una
loss ristretta a queste lettere modella i loro logit relativi, non la loro probabilità assoluta, e il readout si
comporta quindi come un classificatore a tre classi innestato sul modello linguistico. È questo a definire un System
One: la risposta non passa mai per il decoding, ma viene scelta fra le opzioni ammissibili. Il fenomeno richiama anche
l'osservazione che la probabilità del primo token di risposta e la risposta effettivamente scritta dal modello possono
non coincidere~\cite{wang2024myanswerc}, e che la scelta fra lettere è soggetta a bias di selezione~\cite{zheng2024mcq}:
un readout specializzato e calibrato deve assorbire entrambi gli effetti. Un esperimento di controllo che, con gli
stessi pesi e lo stesso passaggio, prende al posto del readout ristretto il token che la generazione emetterebbe ne
quantifica la differenza (Sezione~\ref{sec:results-readout}).

% Marker " Verdict:" = 3 tokens with the Qwen2.5 tokenizer (" Ver", "dict", ":"), verified 2026-10-07.
% Prefill / decode speeds: Phase 0 measurements (Project/docs/phase0_report.md). Readout mass: thesis_notes.md (B2 check).

\section{Obiettivo di addestramento}
\label{sec:method-training}

GroundOne si ottiene da Qwen2.5-0.5B-Instruct e Qwen2.5-3B-Instruct~\cite{qwen2024qwen25} tramite low-rank adaptation
(LoRA)~\cite{hu2022lora}. Gli adapter, di rango 16, con fattore di scala 2 e dropout 0.05, sono applicati a tutti i
layer lineari di tutti i blocchi transformer, poiché coprire ogni layer conta più di un rango
elevato~\cite{dettmers2023qlora}. I pesi del modello base restano congelati in precisione a 8 bit. Un esempio di
addestramento è una risposta con tutte le sue frasi valutate, e la loss viene calcolata su ogni marker della risposta.

La loss è una cross-entropy calcolata sulle sole tre lettere: il prompt non viene mai supervisionato e il segnale di
addestramento proviene interamente dai verdetti. Poiché circa l'88\,\% delle frasi del training set è supportato, le
classi sono pesate in modo che ciascuna contribuisca in egual misura alla loss: con la distribuzione del training set
si ottengono pesi di circa 0.38 per le frasi supportate, 7.5 per quelle contraddette e 4.4 per quelle non menzionate.
Senza questi pesi un modello può ottenere una loss bassa rispondendo quasi sempre ``supportata''. Ne discende una
conseguenza utile: un modello che avesse appreso soltanto le frequenze delle classi resterebbe a una loss di
$\ln 3 \approx 1.10$, quindi una loss inferiore a questo valore indica che il testo viene effettivamente sfruttato.

L'ottimizzazione usa AdamW senza weight decay e, come nella configurazione predefinita di MLX, senza bias correction
delle stime dei momenti; non si applica gradient clipping. Il learning rate di $10^{-4}$ viene raggiunto con un warm-up
lineare sul primo 3\,\% dei passi e poi portato a zero con un cosine schedule, nell'arco di un'epoca sulle risposte del
training set. Le risposte di lunghezza simile sono raggruppate in micro-batch entro un budget di token comprensivo del
padding, e i gradienti vengono accumulati in modo che ogni passo dell'ottimizzatore veda circa 16 risposte, ovvero
circa 110 decisioni: così ogni passo contiene qualche frase contraddetta nonostante la rarità della classe. Il gradient
checkpointing mantiene l'occupazione di memoria entro i 16\,GB del portatile. Il modello viene valutato periodicamente
sul dev set e si conserva il checkpoint con la balanced accuracy a tre classi più alta, mediata sui tre task. La
Sezione~\ref{sec:setup-models} riporta le impostazioni di ciascun run.

\section{Calibrazione}
\label{sec:method-calibration}

I pesi delle classi addestrano il modello come se le tre classi fossero ugualmente frequenti, per cui le sue
probabilità grezze non rispecchiano la frequenza reale degli errori. Sul dev set si stima quindi una trasformazione
post-hoc: i tre logit vengono divisi per una temperatura $T$ e traslati di un bias per classe,
$\mathrm{softmax}(z / T + b)$, con il bias della classe supportata fissato a zero, e $T$ e $b$ sono scelti minimizzando
la negative log-likelihood delle etichette del dev set~\cite{guo2017calibration}. Il bias svolge il ruolo della
correzione contestuale proposta per i modelli linguistici interrogati tramite prompt~\cite{zhao2021calibrate}, qui
però stimata su dati etichettati. La sola temperatura cambia le probabilità ma non la classe più probabile né
l'ordinamento delle frasi. Il bias, invece, può cambiare la classe più probabile e, con tre classi, può anche
riordinare leggermente $P(\text{supported})$ fra una frase e l'altra. Per questo le metriche di classificazione e
l'AUROC sono calcolate sull'output grezzo, mentre le probabilità calibrate si usano dove una probabilità deve
corrispondere a una frequenza: nelle metriche di calibrazione, nelle soglie e nella selective prediction
(Sezione~\ref{sec:results-calibration}).

La calibrazione è misurata principalmente con l'Expected Calibration Error
(ECE)~\cite{naeini2015ece,guo2017calibration}. Le predizioni vengono ripartite in base alla confidenza in $M$
intervalli di uguale ampiezza, con $M = 15$ in questa tesi. Per ogni intervallo $B_m$ si confronta la confidenza media
$\mathrm{conf}(B_m)$ con la frequenza osservata $\mathrm{acc}(B_m)$, e gli scarti vengono mediati con pesi
proporzionali alla numerosità degli intervalli:
\begin{equation}
  \mathrm{ECE} = \sum_{m=1}^{M} \frac{|B_m|}{n} \, \bigl| \mathrm{acc}(B_m) - \mathrm{conf}(B_m) \bigr| .
\end{equation}
Nella variante binaria la confidenza è $P(\text{supported})$ e la frequenza osservata è la quota di frasi supportate
nell'intervallo; nella variante top-label a tre classi la confidenza è la probabilità massima e la frequenza osservata
è la quota di predizioni corrette. L'ECE non confronta la distribuzione delle classi predette con quella del test set:
verifica se, fra le frasi a cui il modello assegna una certa confidenza, quella confidenza corrisponde alla realtà.
Dipende dalla scelta degli intervalli, e un predittore costante che restituisce sempre la frequenza di base ottiene un
ECE vicino a zero pur non discriminando nulla. Per questi motivi l'ECE è sempre riportato insieme al Brier score, alla
negative log-likelihood e ai reliability diagram, mentre la capacità delle probabilità di separare le decisioni
corrette da quelle errate si misura a parte, con l'AUROC e con l'area sotto la risk--coverage curve.

\paragraph{Dalle probabilità alle etichette restituite.} Chi usa il verificatore vede per ogni frase un'etichetta e una
probabilità, che devono essere coerenti tra loro. La decisione grezza, cioè la classe più probabile secondo l'output non
calibrato, è quella valutata in questa tesi, ma può contraddire la probabilità calibrata mostrata accanto: una frase può
essere etichettata ``non menzionata'' pur avendo una $P(\text{supported})$ calibrata di poco superiore a un mezzo.
L'output restituito ricava quindi ogni elemento dalle probabilità calibrate con un'unica regola. Una frase è etichettata
come supportata se e solo se $P(\text{supported}) \geq \tau$; altrimenti riceve la più probabile fra le due classi
restanti. La policy strict segnala ogni frase non supportata, quella lenient soltanto le contraddette. La soglia $\tau$
è stimata sul dev set in modo da massimizzare la balanced accuracy a tre classi, e l'utente può modificarla a ogni
richiesta per scegliere un diverso punto operativo senza riaddestrare il modello.

\section{Software}
\label{sec:method-software}

GroundOne è implementato come pacchetto Python installabile basato su MLX~\cite{mlx2023,mlxlm}, il framework per array
pensato per Apple silicon, e tutte le fasi di addestramento e di inferenza descritte finora vengono eseguite al suo
interno sullo stesso portatile. Il pacchetto comprende il costruttore del dataset con i relativi controlli, il ciclo di
addestramento, la calibrazione, la valutazione sui test set congelati, le baseline e un meccanismo di protezione della
memoria che esegue i processi lunghi uno alla volta e interrompe un processo prima che il sistema esaurisca la memoria.
Il verificatore è esposto tramite un'interfaccia a riga di comando e un server HTTP; il server offre un endpoint
nativo e un endpoint compatibile con l'interfaccia usata da altri modelli System One, così che i client esistenti
possano adottare GroundOne cambiando soltanto l'indirizzo. L'endpoint nativo riceve un documento e una risposta,
segmenta la risposta con la stessa regola usata per costruire i dati, esegue il singolo forward pass della
Sezione~\ref{sec:method-readout} e riporta per ogni frase se è stata valutata, le tre probabilità calibrate,
l'etichetta restituita, i flag delle policy strict e lenient, nonché la regola di decisione e la soglia che li hanno
prodotti; indica inoltre il numero di token in ingresso, il numero di token generati, sempre pari a zero, e il tempo
impiegato. L'endpoint compatibile risponde a domande chiuse generiche su un contesto condiviso, con il formato a una
domanda per frase. Tutto il calcolo del modello avviene su un unico worker thread di lunga durata, perché il contesto
Metal della GPU non deve spostarsi da un thread all'altro, e nulla esce dalla macchina: modelli, dati e richieste
restano in locale.
